\documentclass[11pt]{article}
\usepackage[margin=1in]{geometry}
\usepackage{amsmath,amssymb,amsthm}
\usepackage{xurl}
\usepackage[hidelinks]{hyperref}
\newtheorem{theorem}{Theorem}
\newtheorem{lemma}[theorem]{Lemma}
\newtheorem{corollary}[theorem]{Corollary}
\theoremstyle{remark}
\newtheorem{remark}[theorem]{Remark}
\setlength{\emergencystretch}{3em}

\title{The normalized number of zero signed sums $\pm1\pm2\pm\cdots\pm n$ increases from $n=16$}
\author{Yitzchak Shmalo$^{\,a}$\thanks{Email address:
\texttt{yitzchak.shmalo@gmail.com} (Yitzchak Shmalo).}\\
{\small $^{a}$Einstein Institute of Mathematics, The Hebrew University of
Jerusalem, Givat Ram,}\\
{\small Jerusalem, Israel}}
\date{November 2026}

\begin{document}
\maketitle

\begin{abstract}
Let $S_n=\epsilon_1+2\epsilon_2+\cdots+n\epsilon_n$ with independent fair signs $\epsilon_j$, and let $Q_n=1^2+2^2+\cdots+n^2$ be its variance. On the indices $n\equiv0,3\pmod 4$, where $S_n$ can vanish, the normalized number of zero signed sums $\sqrt{Q_n}\,\Pr(S_n=0)$ tends to $\sqrt{2/\pi}$ by a theorem of Sullivan. We prove that this sequence increases strictly from $n=16$ on, and that the only decreases between consecutive such indices are $3\to4$, $8\to11$ and $15\to16$. When every summand $j\epsilon_j$ is replaced by $m$ independent copies, the corresponding quantity, which also tends to $\sqrt{2/\pi}$, has exactly four decreases for $m=2$, one for $m=3$, and none for $m\ge4$. The proof compares the trigonometric integrals of the sums at consecutive indices on the range where the integrands are nonnegative, using that the vector $(j^2/Q_n)_{j\le n}$ extended by a zero majorizes the corresponding vector for $n+1$. The remaining ranges are bounded separately. The finitely many steps with $mn<320$ are decided by an exact computation.
\end{abstract}

\noindent\textbf{Keywords:} signed sums, trigonometric integrals, majorization, Karamata's inequality, ordered partitions of multisets.

\noindent\textbf{Mathematics Subject Classification (2020):} 05A16, 60F05, 11B75.

\section{Introduction}

The number $Z(n)$ of choices of signs with $\pm1\pm2\pm\cdots\pm n=0$ (sequence A063865 of~\cite{OEIS}) is nonzero exactly when $n(n+1)/2$ is even, that is, when $n\equiv0$ or $3\pmod 4$, and is then the coefficient of $x^{n(n+1)/4}$ in $(1+x)(1+x^2)\cdots(1+x^n)$. Andrica and Tomescu~\cite{AT02} wrote $Z(n)$ as a trigonometric integral, proved $Z(n)\ge6Z(n-4)$ for $n\ge8$~\cite[Theorem~2.3]{AT02}, and conjectured~\cite[Conjecture~2.6]{AT02} that $Z(n)\sim\sqrt{6/\pi}\,2^nn^{-3/2}$. Sullivan~\cite{Sul13} proved this by the Laplace method, and notes that Hwang had suggested a similar approach in his review~\cite{Hwa03} of~\cite{AT02}. Baker and Wagner~\cite{BW16} obtained the second term of the expansion for a wide class of weight sequences.

Let $S_n=\epsilon_1+2\epsilon_2+\cdots+n\epsilon_n$ with independent fair signs $\epsilon_j$, and let $Q_n=1^2+2^2+\cdots+n^2$ be its variance. Then $\Pr(S_n=0)=Z(n)/2^n$, and Sullivan's theorem reads
\[
\sqrt{Q_n}\,\Pr(S_n=0)\to\sqrt{2/\pi}\qquad(n\equiv0,3\bmod 4).
\]
TheoremDB~\cite{TDB} lists as open problem P2538 the question whether this sequence is strictly increasing along $n\equiv0,3\pmod4$ once $n\ge16$.

More generally, let
\[
S_{m,n}=\sum_{j=1}^n j\sum_{i=1}^m\epsilon_{ij},
\]
with independent fair signs $\epsilon_{ij}$, so that each summand $j\epsilon_j$ is replaced by $m$ independent copies; we call $m\ge1$ the multiplicity. Its variance is $mQ_n$, and $Q_n=n(n+1)(2n+1)/6$. All values of $S_{m,n}$ have the parity of $mn(n+1)/2$; call $n\ge1$ \emph{admissible} for $m$ if this number is even, and put
\[
A_{m,n}=\sqrt{mQ_n}\,\Pr(S_{m,n}=0)\qquad(n\ \text{admissible}).
\]
We call $A_{m,n}$ the \emph{normalized number of zero signed sums}. For odd $m$ the admissible indices are those with $n\equiv0,3\pmod4$, and for even $m$ every $n\ge1$ is admissible; in both cases consecutive admissible indices differ by at most $3$. A \emph{step} is a pair $n<a$ of consecutive admissible indices, and it is a \emph{decrease} if $A_{m,a}<A_{m,n}$. For $m=1$ we have $S_{1,n}=S_n$, so $A_{1,n}=\sqrt{Q_n}\,\Pr(S_n=0)$.

The number $C_{m,n}=2^{mn}\Pr(S_{m,n}=0)$ is the middle coefficient of $\prod_{j=1}^n(1+x^j)^m$. It counts the ordered $2$-partitions of equal sum of the multiset with $m$ labelled copies of each of $1,\dots,n$, that is, the ordered pairs of disjoint parts with equal sums whose union is this multiset, and $C_{1,n}=Z(n)$.

Bagdasar and Andrica~\cite{BA17} gave recurrences, generating functions and integral formulas for the numbers of ordered $2$-partitions of equal sum of multisets with repeated elements, Andrica and Bagdasar~\cite{AB18} did the same for $3$-partitions, and in~\cite{AB21} they derived the generating function and an integral formula for ordered partitions of a multiset into any number of parts of equal sum. Both~\cite{BA17} and~\cite{AB21} formulate conjectures on the asymptotic behaviour of these numbers, and He, Lv and Zhou~\cite{HLZ23}, confirming conjectures of Andrica and Bagdasar, proved that $C_{m,n}\sim2^{mn}n^{-3/2}\sqrt{6/(\pi m)}$ along admissible $n$, together with the corresponding formulas for more parts. The case $m=2$ is older. The number of subsets of $\{-n,\dots,n\}$ with sum $0$ is $2C_{2,n}$: each pair $\{-j,j\}$ contributes $j(\epsilon+\epsilon')/2$ with two fair signs, and the element $0$ doubles the count. Van Lint~\cite{vL67} found the asymptotic number of these subsets, and Entringer~\cite{Ent68} and Clark~\cite{Cla00} extended it to subsets with other prescribed sums. Since $Q_n\sim n^3/3$, the formula of He, Lv and Zhou says that $A_{m,n}\to\sqrt{2/\pi}$ for every fixed $m$. An asymptotic formula, even with its second term, does not decide the sign of $A_{m,a}-A_{m,n}$ at a given step.

For $m=1$ the theorem of Baker and Wagner~\cite[Theorem~1 and Section~2.2.1]{BW16} gives
\[
A_{1,n}=\sqrt{2/\pi}\,\Bigl(1-\frac{r_n}{4}\Bigr)+O\bigl(n^{-2}\bigr),\qquad
r_n=\frac{1}{Q_n^2}\sum_{j\le n}j^4\sim\frac{9}{5n},
\]
so for large $n$ the sequence $A_{1,n}$ approaches $\sqrt{2/\pi}$ from below. For a step $n<a$ the two-term values $\sqrt{2/\pi}\,(1-r_n/4)$ and $\sqrt{2/\pi}\,(1-r_a/4)$ differ by an amount of order $n^{-2}$, which the error term does not control.

\begin{theorem}\label{thm:main}
Let $m\ge1$, and let $n<a$ be a step. Then $A_{m,a}\ne A_{m,n}$, and the step is a decrease exactly in the following eight cases:
\begin{align*}
m=1&:\quad (n,a)\in\{(3,4),(8,11),(15,16)\},\\
m=2&:\quad (n,a)\in\{(3,4),(4,5),(5,6),(6,7)\},\\
m=3&:\quad (n,a)=(3,4).
\end{align*}
\end{theorem}

\begin{corollary}\label{cor:one}
Along $n\equiv0,3\pmod 4$, the sequence $A_{1,n}=\sqrt{Q_n}\,\Pr(S_n=0)$ is strictly increasing for $n\ge16$, and the step from $15$ to $16$ is a decrease. For every $m\ge4$ the sequence $A_{m,n}$ is strictly increasing along all admissible $n$.
\end{corollary}

For admissible $n$,
\begin{equation}\label{eq:A}
A_{m,n}=\frac2\pi\int_0^{B_n}f_n(x)\,dx,\qquad
f_n(x)=\prod_{j=1}^n\cos^m\Bigl(\frac{jx}{\sqrt{mQ_n}}\Bigr),\qquad B_n=\frac\pi2\sqrt{mQ_n}.
\end{equation}
We compare these integrals at the two indices of a step. Every factor of $f_n$ is nonnegative on $[0,b_n]$, where $b_n=\pi\sqrt{mQ_n}/(2n)$. The vector $(j^2/Q_n)_{j\le n}$, extended by a zero, majorizes $(j^2/Q_{n+1})_{j\le n+1}$: both vectors have sum $1$, and the partial sums of the first, in decreasing order, dominate those of the second. Moreover $-\log\cos z$ is a power series in $z^2$ with positive coefficients. Karamata's inequality~\cite[Chapter~3, Proposition~C.1]{MOA11}, applied to $y\mapsto y^l$ for each power $z^{2l}$ of this series, gives for $N>n$
\[
f_N(x)\ \ge\ f_n(x)\exp\Bigl(\frac{x^4}{12m}\,(r_n-r_{n+1})\Bigr)\quad(0\le x\le b_n),\qquad r_n-r_{n+1}>\frac{1}{4n^2}.
\]
On $[1,2]$ this gives a lower bound for $\int_0^{b_n}(f_a-f_n)\,dx$ (Lemma~\ref{lem:gain}), and the parts of~\eqref{eq:A} beyond $b_n$ and $b_a$ are bounded as in Sullivan~\cite[Lemma~2, Case~3]{Sul13} (Lemma~\ref{lem:tail}). For $mn\ge320$ the lower bound of Lemma~\ref{lem:gain} exceeds the sum of the two bounds of Lemma~\ref{lem:tail}, and $f_a\ge0$ on $[b_n,b_a]$, so $A_{m,a}>A_{m,n}$. For each of the $1268$ steps with $mn<320$ the sign of $A_{m,a}-A_{m,n}$ is the sign of an explicit integer, which is evaluated exactly.

\section{Proof}\label{sec:proof}

Fix $m$ and write $\varphi_n(t)=\prod_{j=1}^n\cos^m(jt)$, the characteristic function of $S_{m,n}$. Replacing $t$ by $t+\pi$ multiplies $\varphi_n$ by $(-1)^{mn(n+1)/2}$, so for admissible $n$ the even function $\varphi_n$ has period $\pi$, and Fourier inversion gives $\Pr(S_{m,n}=0)=\frac{2}{\pi}\int_0^{\pi/2}\varphi_n(t)\,dt$, which is the trigonometric integral of Andrica and Tomescu~\cite[Section~1]{AT02} with every multiplicity equal to $m$. The substitution $x=t\sqrt{mQ_n}$ gives~\eqref{eq:A}, since $f_n(x)=\varphi_n\bigl(x/\sqrt{mQ_n}\bigr)$.
The functions $f_n$ and the numbers $B_n$, $b_n$ and $r_n$ are defined by the same formulas for every $n\ge1$, admissible or not. For $0\le x\le b_n$ every argument $jx/\sqrt{mQ_n}$ with $j\le n$ lies in $[0,\pi/2]$, so $f_n\ge0$ there. Since $Q_n/n^2=(2n+3+1/n)/6$ increases with $n$, so does $b_n$.

\begin{lemma}\label{lem:pointwise}
For $N>n\ge1$ and $0\le x\le b_n$,
\[
f_N(x)\ge f_n(x)\exp\Bigl(\frac{x^4}{12m}(r_n-r_{n+1})\Bigr),
\]
and $r_n-r_{n+1}>1/(4n^2)$.
\end{lemma}

\begin{proof}
On $|z|<\pi/2$ we have $-\log\cos z=\sum_{l\ge1}c_lz^{2l}$ with $c_1=1/2$, $c_2=1/12$ and every $c_l>0$. Indeed, the series is the integral of the Taylor series of $\tan z$, whose nonzero coefficients are positive because $\tan z$ solves $y'=1+y^2$, $y(0)=0$, and whose radius of convergence is $\pi/2$. Let $k\ge n$ and $0\le x<b_k$. Then
\[
-\log f_k(x)=m\sum_{l\ge1}c_l\Bigl(\frac{x^2}{m}\Bigr)^lP_l(k),\qquad P_l(k)=\sum_{j=1}^k\Bigl(\frac{j^2}{Q_k}\Bigr)^l,
\]
and the same holds with $k+1$ in place of $k$ because $b_k<b_{k+1}$.

We claim $P_l(k)\ge P_l(k+1)$ for every $l\ge1$, with equality for $l=1$ since $P_1\equiv1$. Let $u\in\mathbb R^{k+1}$ list $k^2/Q_k,\dots,1/Q_k,0$ and $w$ list $(k+1)^2/Q_{k+1},\dots,1/Q_{k+1}$, both in decreasing order and both with sum $1$. For $i\le k$ the ratio $w_i/u_i=\bigl((j+1)/j\bigr)^2Q_k/Q_{k+1}$ with $j=k+1-i$ increases with $i$, and $u_{k+1}-w_{k+1}<0$. Hence $u_i-w_i$ changes sign at most once, and since both vectors have sum $1$ its signs run from positive to negative; the partial sums of $u-w$ rise and then fall from $0$ back to $0$, and so they are nonnegative. Thus $u$ majorizes $w$, and Karamata's inequality~\cite[Chapter~3, Proposition~C.1]{MOA11} for the convex functions $y\mapsto y^l$ on $[0,\infty)$ gives the claim. Subtracting the two series and keeping only the term $l=2$, where $mc_2(x^2/m)^2=x^4/(12m)$, gives for $0\le x<b_k$
\[
f_{k+1}(x)\ge f_k(x)\exp\Bigl(\frac{x^4}{12m}\bigl(P_2(k)-P_2(k+1)\bigr)\Bigr),
\]
and by continuity this also holds at $x=b_k$, where $f_k$ vanishes.

For $x\le b_n\le b_k$ the functions $f_k$, $n\le k\le N$, are nonnegative and each exponential factor is at least $1$, so we may chain these inequalities over $k=n,\dots,N-1$, keeping the exponential factor only for $k=n$. Since $P_2(k)=r_k$, this proves the inequality.

Finally $\sum_{j\le n}j^4=n(n+1)(2n+1)(3n^2+3n-1)/30$ gives
\[
r_n=\frac{6(3n^2+3n-1)}{5n(n+1)(2n+1)},\qquad r_n-r_{n+1}=\frac{36(n^2+2n-1)}{5n(n+2)(2n+1)(2n+3)},
\]
and $144n^2(n^2+2n-1)-5n(n+2)(2n+1)(2n+3)=n(124n^3+208n^2-239n-30)>0$ for $n\ge1$.
\end{proof}

\begin{lemma}\label{lem:gain}
If $mn\ge320$ and $N>n$, then $\int_0^{b_n}(f_N-f_n)\,dx\ge e^{-4}/(48mn^2)$.
\end{lemma}

\begin{proof}
The integrand is nonnegative by Lemma~\ref{lem:pointwise}. Since $Q_n\ge n^3/3$, for $1\le x\le2$ and $j\le n$ we have $jx/\sqrt{mQ_n}\le2\sqrt{3/(mn)}<1/2$, and $b_n\ge\frac\pi2\sqrt{mn/3}>2$. On $[0,1/2]$, $\tan s\le2s$ because $\sec^2s\le2$, so $-\log\cos z\le z^2$, and therefore $f_n(x)\ge\exp(-\sum_j j^2x^2/Q_n)=e^{-x^2}\ge e^{-4}$ on $[1,2]$. By Lemma~\ref{lem:pointwise} and $e^y-1\ge y$, on $[1,2]$
\[
f_N-f_n\ge f_n\,\frac{x^4}{12m}\cdot\frac{1}{4n^2}\ge\frac{e^{-4}}{48mn^2}.
\]
Since $[1,2]\subset[0,b_n]$ has length $1$ and the integrand is nonnegative on the rest of $[0,b_n]$, the lemma follows.
\end{proof}

The next estimate is Sullivan's~\cite[Lemma~2, Case~3]{Sul13}, extended to multiplicity $m$ and to the range $t\ge\pi/(2k)$, with the bound $1-s\le e^{-s}$ in place of the inequality of arithmetic and geometric means.

\begin{lemma}\label{lem:tail}
For every $k\ge1$, $\int_{b_k}^{B_k}|f_k(x)|\,dx\le\frac\pi2\sqrt{mQ_k}\,e^{-mk/8}$.
\end{lemma}

\begin{proof}
With $x=t\sqrt{mQ_k}$ the integral is $\sqrt{mQ_k}\int|\varphi_k(t)|\,dt$ over $\pi/(2k)\le t\le\pi/2$, where $\sin t\ge2t/\pi\ge1/k$. Summing a geometric series,
\[
\sum_{j=1}^k\sin^2(jt)=\frac k2+\frac14-\frac{\sin((2k+1)t)}{4\sin t}\ge\frac k2+\frac14-\frac k4\ge\frac k4 .
\]
Since $\cos^2\theta=1-\sin^2\theta\le e^{-\sin^2\theta}$, we get $|\varphi_k(t)|\le\exp\bigl(-\frac m2\sum_j\sin^2(jt)\bigr)\le e^{-mk/8}$, and the range of $t$ has length less than $\pi/2$.
\end{proof}

\begin{proof}[Proof of Theorem~\ref{thm:main} for $mn\ge320$]
Let $n<a$ be a step with $mn\ge320$; then $a\le n+3\le4n$. By~\eqref{eq:A}, splitting both integrals at $b_n$ and the first again at $b_a>b_n$,
\[
\frac\pi2\bigl(A_{m,a}-A_{m,n}\bigr)=\int_0^{b_n}(f_a-f_n)\,dx+\int_{b_n}^{b_a}f_a\,dx+\int_{b_a}^{B_a}f_a\,dx-\int_{b_n}^{B_n}f_n\,dx .
\]
The second integral is nonnegative. By Lemmas~\ref{lem:gain} and~\ref{lem:tail}, with $Q_k\le k^3$, $a\le4n$ and $e^{-ma/8}\le e^{-mn/8}$,
\[
\frac\pi2\bigl(A_{m,a}-A_{m,n}\bigr)\ge\frac{e^{-4}}{48mn^2}-\frac\pi2\sqrt m\,\bigl(n^{3/2}+8n^{3/2}\bigr)e^{-mn/8}.
\]
The right side is positive as soon as $216\pi e^4m^{3/2}n^{7/2}e^{-mn/8}<1$, which follows from $216\cdot4\cdot81\cdot v^4e^{-v/8}<1$ with $v=mn$ (using $\pi<4$, $e^4<81$ and $m^{3/2}n^{7/2}\le(mn)^4$). The function $v^4e^{-v/8}$ decreases for $v\ge32$, and at $v=320$ the inequality reads $216\cdot4\cdot81\cdot320^4<e^{40}$, which holds because $216\cdot4\cdot81\cdot320^4<7.4\cdot10^{14}$ and $e^{40}>(8/3)^{40}>10^{17}$.
\end{proof}

\begin{proof}[Proof of Theorem~\ref{thm:main} for $mn<320$]
Here $1\le m\le319$ and $n\le319/m$. Since $C_{m,n}\ge0$ and $A_{m,n}=\sqrt{mQ_n}\,C_{m,n}/2^{mn}$, the difference $A_{m,a}-A_{m,n}$ has the sign of the integer
\[
Q_aC_{m,a}^2-4^{m(a-n)}Q_nC_{m,n}^2 .
\]
For the $1268$ steps with $mn<320$ these integers are computed exactly from the coefficients $C_{m,n}$ of $\prod_{j\le n}(1+x^j)^m$; none is zero, and exactly the eight listed in the theorem are negative. The integers are computed by the program \texttt{check\_polynomial.py} listed under Code and data availability.
\end{proof}

\begin{remark}\label{rem:values}
Rounded to four places, $A_{1,16}=0.7755$, $A_{1,100}=0.7943$ and $A_{1,300}=0.7967$, and $\sqrt{2/\pi}=0.7979$. For $n=100$ and $n=300$, $|A_{1,n}-\sqrt{2/\pi}\,(1-r_n/4)|<10^{-5}$.
\end{remark}

\section*{Acknowledgments and funding}
The research presented in this paper was supported by the European Research Council (ERC) under the European Union's Horizon Europe research and innovation programme (grant agreement No.~101041711), and by the Simons Foundation as part of the Collaboration on the Mathematical and Scientific Foundations of Deep Learning. It was also supported by the Israel Science Foundation (grants 2258/19 and 4101/25), by Heights Labs and by Convex Nexus Capital.

This work was done by multiple AI agents, Claude (Anthropic) and Codex (OpenAI), connected to Math Brain~\cite{MathBrainRepo} and working together inside the Brain framework~\cite{BrainFramework}. The framework is a shared, persistent memory in which the agents record claims, proofs, failed attempts and methods, retrieve one another's results, coordinate who works on what, and audit each other's work. The method is described in~\cite{MathBrainPaper}.

\section*{Declaration of competing interest}
The author holds equity in Heights Labs and in Convex Nexus Capital and has a management role at Convex Nexus Capital. Both are among the funders named above.

\section*{Code and data availability}
The \LaTeX\ source and two programs with their outputs are available at \url{https://github.com/yspennstate/signed-sums-zero-monotone}. The program \texttt{check\_polynomial.py} decides the steps with $mn<320$ by multiplying out $\prod_{j\le n}(1+x^j)^m$ and prints the values in Remark~\ref{rem:values}. The program \texttt{check\_signed\_walk.py} decides the same steps by counting choices of signs.

\section*{Declaration of generative AI and AI-assisted technologies}
This paper is part of a broader research project on proving open problems with AI systems. Codex (OpenAI) wrote the proof of Theorem~\ref{thm:main} and the first computation of the steps with $mn<320$. Claude (Anthropic) checked the proof and wrote the two programs of the repository and the text.

The question is problem P2538 of TheoremDB~\cite{TDB}. Earlier results on it are due to Andrica and Tomescu~\cite{AT02}, Sullivan~\cite{Sul13}, Baker and Wagner~\cite{BW16}, van Lint~\cite{vL67}, Bagdasar and Andrica~\cite{BA17,AB21} and He, Lv and Zhou~\cite{HLZ23}, and are described in the introduction. The author directed the work and is responsible for the content of the paper and for decisions concerning its dissemination.

\begin{thebibliography}{99}

\bibitem{AB18}
D.~Andrica and O.~Bagdasar.
\newblock Some remarks on 3-partitions of multisets.
\newblock \emph{Electron. Notes Discrete Math.}, 70:1--8, 2018.
\newblock \url{https://doi.org/10.1016/j.endm.2018.11.001}.

\bibitem{AB21}
D.~Andrica and O.~Bagdasar.
\newblock On $k$-partitions of multisets with equal sums.
\newblock \emph{Ramanujan J.}, 55(2):421--435, 2021.
\newblock \url{https://doi.org/10.1007/s11139-021-00418-7}.

\bibitem{AT02}
D.~Andrica and I.~Tomescu.
\newblock On an integer sequence related to a product of trigonometric functions, and its combinatorial relevance.
\newblock \emph{J. Integer Seq.}, 5(2):Article 02.2.4, 2002.
\newblock \url{https://cs.uwaterloo.ca/journals/JIS/VOL5/Tomescu/tomescu4.html}.

\bibitem{BA17}
O.~Bagdasar and D.~Andrica.
\newblock New results and conjectures on 2-partitions of multisets.
\newblock In \emph{2017 7th International Conference on Modeling, Simulation, and Applied Optimization (ICMSAO)}, pages 1--5, 2017.
\newblock \url{https://doi.org/10.1109/ICMSAO.2017.7934928}.

\bibitem{BW16}
L.~Baker and S.~Wagner.
\newblock Erd\H{o}s--Sur\'anyi sequences and trigonometric integrals.
\newblock \emph{Publ. Math. Debrecen}, 89(4):469--482, 2016.
\newblock \url{https://doi.org/10.5486/PMD.2016.7452}; arXiv:1506.04555.

\bibitem{Cla00}
L.~Clark.
\newblock On the representation of $m$ as $\sum_{k=-n}^{n}\epsilon_k k$.
\newblock \emph{Int. J. Math. Math. Sci.}, 23(1):77--80, 2000.
\newblock \url{https://doi.org/10.1155/S0161171200001885}.

\bibitem{Ent68}
R.~C. Entringer.
\newblock Representation of $m$ as $\sum_{k=-n}^{n}\epsilon_k k$.
\newblock \emph{Canad. Math. Bull.}, 11:289--293, 1968.
\newblock \url{https://doi.org/10.4153/CMB-1968-036-3}.

\bibitem{HLZ23}
X.~He, J.-B. Lv, and N.~H. Zhou.
\newblock Note on the number of ordered $k$-partitions of multiset $M_d[n]$ with equal sums.
\newblock \emph{Mediterr. J. Math.}, 20(4):Paper No.~228, 2023.
\newblock \url{https://doi.org/10.1007/s00009-023-02435-w}.

\bibitem{Hwa03}
H.-K. Hwang.
\newblock Review of D. Andrica and I. Tomescu, On an integer sequence related to a product of trigonometric functions, and its combinatorial relevance.
\newblock \emph{Mathematical Reviews}, MR1938223 (2003j:05005).

\bibitem{MOA11}
A.~W. Marshall, I.~Olkin and B.~C. Arnold, \emph{Inequalities: Theory of Majorization and Its Applications}, 2nd ed., Springer, New York, 2011.
\newblock \url{https://doi.org/10.1007/978-0-387-68276-1}.

\bibitem{Sul13}
B.~D. Sullivan.
\newblock On a conjecture of Andrica and Tomescu.
\newblock \emph{J. Integer Seq.}, 16:Article 13.3.1, 2013.
\newblock \url{https://cs.uwaterloo.ca/journals/JIS/VOL16/Sullivan/sullivan8.html}; arXiv:1210.8437.

\bibitem{vL67}
J.~H. van Lint.
\newblock Representation of $0$ as $\sum_{k=-N}^{N}\epsilon_k k$.
\newblock \emph{Proc. Amer. Math. Soc.}, 18:182--184, 1967.
\newblock \url{https://doi.org/10.1090/S0002-9939-1967-0205964-8}.

\bibitem{OEIS}
OEIS Foundation Inc.
\newblock The On-Line Encyclopedia of Integer Sequences, sequence A063865.
\newblock \url{https://oeis.org/A063865}, accessed 20 November 2026.

\bibitem{TDB}
TheoremDB, problem P2538.
\newblock \url{https://theoremdb.org/statements/P2538/}, accessed 20 November 2026.

\bibitem{MathBrainRepo}
Y.~Shmalo.
\newblock \emph{Math Brain: public snapshot, code and tests}, 2026.
\newblock \url{https://github.com/yspennstate/math-brain-public}.

\bibitem{MathBrainPaper}
Y.~Shmalo.
\newblock \emph{Math Brain: shared, auditable memory for multi-agent mathematical research}, October 2026.
\newblock \url{https://github.com/yspennstate/math-brain-public/tree/main/papers/math-brain-2026-10-24}.

\bibitem{BrainFramework}
Y.~Shmalo.
\newblock \emph{Brain Framework}, public software repository, 2026.
\newblock \url{https://github.com/yspennstate/brain-framework-public}.

\end{thebibliography}
\end{document}
